% -----------------------------------------------------------
\section{Covenant}
% -----------------------------------------------------------
	\label{COVENANT}
	
% % ----------------------
% \subsection{See also}
% % ----------------------

% ----------------------
\subsection{1828 American Dictionary of the English Language} % *1828
% ----------------------
	\label{COVENANT-1828}
	
	% -----
	\subsubsection{Noun}
	% -----

		\begin{compactitem}
			\item[1] A mutual consent or agreement of two or more persons, to do or to forbear some act or thing; a contract; stipulation. A \textit{covenant} is created by deed in writing, sealed and executed; or it may be implied in the contract.
			\item[2] A writing containing the terms of agreement or contract between parties; or the clause of agreement in a deed containing the \textit{covenant}.
			\item[3] In theology, the \textit{covenant} of works, is that implied in the commands, prohibitions, and promises of God; the promise of God to man, that man's perfect obedience should entitle him to happiness. This do, and live; that do, and die. The \textit{covenant} of redemption, is the mutual agreement between the Father and Son, respecting the redemption of sinners by Christ. The \textit{covenant} of grace, is that by which God engages to bestow salvation on man, upon the condition that man shall believe in Christ and yield obedience to the terms of the gospel.
			\item[4] In church affairs, a solemn agreement between the members of a church, that they will walk together according to the precepts of the gospel, in brotherly affection.
		\end{compactitem}
		
	% -----
	\subsubsection{Intransitive verb}
	% -----
	
		\begin{compactitem}
			\item To enter into a formal agreement; to stipulate; to bind ones self by contract. A \textit{covenant}s with B to convey to him a certain estate. When the terms are expressed ti has for before the thing or price.
		\end{compactitem}
		
	% -----
	\subsubsection{Transitive verb}
	% -----
	
		\begin{compactitem}
			\item To grant or promise by \textit{covenant}.
		\end{compactitem}

% ----------------------
\subsection{Oxford English Dictionary} % *OED
% ----------------------
	\label{COVENANT-OED}
	
	% -----
	\subsubsection{Noun}
	% -----

		\begin{compactitem}
			\item[1a] A mutual agreement between two or more persons to do or refrain from doing certain acts; a compact, contract, bargain; sometimes, the undertaking, pledge, or promise of one of the parties. Phrases, \textbf{\textit{to make} or \textit{enter into a covenant; to hold, keep, break covenant}} (No longer in ordinary use, except when coloured by legal or theological associations.)
			\item[1b] \textbf{\textit{to, on, upon, in, at (a} or \textit{the) covenant}}: on a mutual stipulation, or understanding; on the condition \textit{that. Obsolete.}
			\item[2] A promise made to oneself, a solemn personal resolve, a vow. \textit{Obsolete}.
			\item[3] Each of the points or terms of an agreement. \textit{Obsolete} except as in 4b.
			\item[4a] A formal agreement, convention, or promise of legal validity; esp. in English Law, a promise or contract under seal. (The English equivalent of Latin conventio as technically used from the Norman Conquest onwards.)
			\item[4b] \textit{esp}. A particular clause of agreement contained in a deed; e.g. the ordinary covenants to pay rent, etc. in a lease.
			\item[5] The matter agreed upon between two parties, or undertaken or promised by either; hence, covenanted duty, service, wages, rent, etc. \textit{Obsolete}.
			\item[6] Pledge, security. \textit{Obsolete}. \textit{rare}.
			\item[7a] \textit{Scripture}. Applied esp. to an engagement entered into by the Divine Being with some other being or persons.
				\begin{compactitem}
					\item (The Hebrew word \hyperref[BERIT]{\hel{b*:riyt}} is also the ordinary term for a contract, agreement, alliance, or league between men. It is constantly rendered in the Septuagint by \hyperref[DIATHEKE]{\grl{διαθήκη}} `disposition, distribution, arrangement', which occurs in Aristophanes in the sense `convention, arrangement between parties', but usually in classical Greek meant `disposition by will, testament'. Accordingly, the Old Latin translation of the Bible (Itala) appears to have uniformly rendered \hyperref[DIATHEKE]{\grl{διαθήκη}} by \textit{testamentum}, while Jerome translated the Hebrew by \textit{foedus} and \textit{pactum} indifferently. Hence, in the Vulgate, the Old Testament has the old rendering \textit{testamentum} in the (Gallican) Psalter, but Jerome's renderings \textit{foedus}, \textit{pactum} elsewhere; the New Testament has always \textit{testamentum}. In English Wyclif strictly followed the Vulgate, rendering \textit{foedus}, \textit{pactum}, by \textit{boond}, \textit{covenaunt}, rather indiscriminately, \textit{testamentum} in the Psalter and New Testament always by \textit{testament}. So the versions of Rheims and Douay. The 16th cent. English versions at length used \textit{covenant} entirely in Old Testament (including the Psalter), and Tyndale introduced it into 6 places in the New Testament. These the Geneva extended to 23, and the Bible of 1611 to 22 (in 2 of which Genesis had \textit{testament}), leaving \textit{testament} in 14 (in 3 of which Genesis had \textit{covenant}). The Revised Version of 1881 has substituted \textit{covenant} in 12 of these, leaving \textit{testament} in 2 only (Hebrew ix. 16, 17).)
					\item Thus \hyperref[BERIT]{\hel{b*:riyt}}, \hyperref[DIATHEKE]{\grl{διαθήκη}}, \textit{fœdus} (\textit{pactum}), \textit{covenant} are applied to God's engagement with Noah and his posterity, Genesis vi. 18, ix. 9–17; to that made with Abraham and his posterity, Genesis xvii, of which the token was circumcision; to the institution of the Mosaic law, Exodus xxiv. 7, 8, and to that law or its observance itself, whence the expressions \textit{book of the covenant} (i.e. of the law), \textit{ark of the covenant}, \textit{blood of the covenant} (i.e. of beasts ritually sacrificed), \textit{land of the covenant} (= promised land, Canaan). The covenant with the Israelites, in its various phases, is commonly called the \textit{Old Covenant}, in contrast to which the prophets made promise of a \textit{new covenant}, Jeremiah xxxi. 31; and this name \gr{καινὴ} \hyperref[DIATHEKE]{\grl{διαθήκη}} (\textit{New Covenant} or \textit{testament}) was, according to St. Luke xxii. 20, applied by Jesus to the new relation to man which God had established in Him. In this sense it is also used by St. Paul and the writer of the Epistle to the Hebrews, who contrast these two covenants (Galatians iv. 24, Hebrew viii. 13, ix. 15, etc.), also called by commentators the \textit{Temporal} and the \textit{Eternal Covenant} (cf. Hebrew xiii. 20).]
				\end{compactitem}
			\item[7b] Hence \textit{covenant} is sometimes used = Dispensation.
			\item[7c] The two divisions of the Scriptures, belonging to the Mosaic and Christian dispensations respectively, are sometimes called the \textbf{\textit{Books of the Old and the New Covenant}}, instead of the usual form Old and New Testament (Greek \gr{παλαιὰ} and \gr{καινὴ} \hyperref[DIATHEKE]{\grl{διαθήκη}}).
			\item[7d] \textbf{\textit{(Greater) Book of the Covenant, Little Book of the Covenant}}: names given by Old Testament critics to certain portions of the Book of Exodus, viz. ch. xx. 22--xxiii, and ch. xxxiv. 11--26 respectively.
			\item[8] \textit{Theology}.
			\item[8a] \textbf{\textit{Covenant of Works, Covenant of Grace}}: the two relations which are represented as subsisting between God and man, before and since the Fall.
			\item[8b] Applied to the engagement with God which is entered into by believers at their baptism, or admission into the visible church.
			\item[9] \textit{Ecclesiastical}.
			\item[9a] \textit{Scottish History}. The name given to certain bonds of agreement signed by the Scottish Presbyterians for the defence and furtherance of their religion and ecclesiastical polity.
			\item[9b] \textit{church covenant}. See ``church'' noun and adjective.
		\end{compactitem}

% ----------------------
\subsection{Scriptural Usage} % *SCRIPTURES
% ----------------------
	\label{COVENANT-SCRIPTURES}

	% -----
	\subsubsection{Old Testament} % *OT
	% -----
		\label{COVENANT-OT}
	
		% % --
		% \paragraph{KJV}
		% % --
			% \label{COVENANT-OT-KJV}
	
			% \begin{tabular}{ll}
				% Total occurrences in the OT & 
			% \end{tabular}
			
			% \begin{compactdesc}
				% \item[]
			% \end{compactdesc}
			
		% --
		\paragraph{Hebrew Bible}
		% --
			\label{COVENANT-HEBREW}
		
			%
			\subparagraph{\texorpdfstring{\he{b*:riyt}}{}}
			%
				\label{BERIT} % whatever is in step bible without periods
				
				See the \hyperref[COVENANT-reference]{reference information} below.
			
				\begin{compactdesc}
					\item[HALOT 157a, PDF 214] \textbf{agreement, covenant}
						\begin{compactitem}
							\item[A] between persons
							\item[B] \textbf{contract, covenant} with---
							\item[B1] God with animals
							\item[B2] with the stones of the fields, with death, through a sacrifice to a god of death
							\item[C] a \textbf{covenant between God and mankind}
							\item[Ca] a covenant which is made
							\item[Ca1] God made a covenant
							\item[Ca2] God made a covenant for the benefit of
							\item[Ca3] God established his covenant with
							\item[Ca4] God enters into a covenant with
							\item[Ca5] God makes an everlasting promise
							\item[Ca6] as a representative of the community, a human individual makes a covenant with God
							\item[Ca7] to enter into a covenant with God
							\item[Cb] phrases connected with \he{b*:riyt}...
							\item[Cc] the maintenance of the covenant
							\item[Cd] neglecting, contravening the covenant
							\item[C3] miscellaneous uses
						\end{compactitem}
					\item[BDB 136a, PDF 155] \textbf{covenant}; \textit{pact, compact, covenant}
						\begin{compactitem}
							\item[I] between men
							\item[I1] \textit{treaty, alliance}
							\item[I2] \textit{constitution, ordinance}
							\item[I3] \textit{agreement, pledge}
							\item[I4] \textit{alliance} of friendship
							\item[I5] \textit{alliance} of marriage
							\item[II] between God and man
							\item[II1] \textit{alliance} of friendship
							\item[II2] \textit{covenant}, as a divine constitution or ordinance with signs and pledges
							\item[II2a] with Noah
							\item[II2b] with Abraham, Isaac and Jacob
							\item[II2c] with Israel at Sinai = Horeb, with a covenant sacrifice...renewed in plains of Moab...with blessings and curses (\hyperref[DEUTERONOMY-29-20]{Deuteronomy 29:20})...frequently referred to in other books...a divine constitution given to Israel with promises on condition of obedience and penalties for disobedience, in the form of tables of the covenant (\hyperref[DEUTERONOMY-9-9]{Deuteronomy 9:9, 11, 15}), inscribed with the ten words, placed in the ark of the covenant
							\item[II2d] with Phinehas
							\item[II2e] with Joshua and Israel
							\item[II2f] with David
							\item[II2g] Jehoiada and the people
							\item[II2h] Hezekiah and the people
							\item[II2i] Josiah and the people
							\item[II2j] Ezra and the people
							\item[II2k] \textit{the prophetic covenant}, a divine promise through a series of prophets to establish a new constitution
							\item[III] Phrases...
						\end{compactitem}
					\item[DCH PDF 739] \textbf{covenant}; \textbf{covenant, agreement}, or \textbf{obligation} between individuals (for example, friends, spouses) or groups, ruler and subjects, deity and individual or people, etc.
						\begin{compactitem}
							\item (The remainder of this entry gives examples of the uses of \he{b*:riyt} as subject, object, construct form, in prepositions, various phrases, and so on.)
						\end{compactitem}
				\end{compactdesc}
				
				% \begin{compactdesc}
					% \item[Some OT uses] \textbf{}
						% \begin{compactdesc}
							% \item[]
						% \end{compactdesc}
				% \end{compactdesc}
			
		% --
		\paragraph{Septuagint}
		% --
			\label{COVENANT-LXX}
		
			%
			\subparagraph{\texorpdfstring{\gr{διαθήκη}}{}}
			%
			
				\begin{compactdesc}
					\item[GELS 150b, PDF 150] \textbf{}
						\begin{compactitem}
							\item[1] \textit{compact, treaty, mutual agreement}: God initiating and mediating in establishment of such between Israel and the animal kingdom...with humans...of a king proposing such with his subjects...
							\item[2a] ``\textit{covenant}'' between God and Israel: but spoken of as God's
							\item[2b] also with a human referent in the genitive
						\end{compactitem}
				\end{compactdesc}
		
	% -----
	\subsubsection{New Testament} % *NT
	% -----
		\label{COVENANT-NT}
	
		% % --
		% \paragraph{KJV}
		% % --
			% \label{COVENANT-NT-KJV}
			
	
			% \begin{tabular}{ll}
				% Total occurrences in the NT & 
			% \end{tabular}
			
			% \begin{compactdesc}
				% \item[]
			% \end{compactdesc}
			
		% --
		\paragraph{Greek New Testament} % *GREEK
		% --
			\label{COVENANT-GREEK}
		
			%
			\subparagraph{\texorpdfstring{\gr{διαθήκη}}{}}
			%
				\label{DIATHEKE}
				
				To understand this work in Greek and in the NT, you have to start with the fact that \gr{διαθήκη} was the word to translate \he{b*:riyt} in the LXX. See BDAG meaning \#2, and also the \hyperref[COVENANT-reference]{reference information} below.
			
				\begin{compactdesc}
					\item[BDAG 202a] \textbf{}
						\begin{compactitem}
							\item[1] \textit{last will and testament}
							\item[2] As a translation of \he{b*:riyt} in LXX, \gr{διαθήκη} retains the component of legal disposition of personal goods while omitting that of the anticipated death of a testator. A Hellenistic reader would experience no confusion, for it was a foregone conclusion that gods were immortal. Hence a \gr{διαθήκη} decreed by God cannot require the death of the testator to make it operative. Nevertheless, another essential characteristic of a testament is retained, namely that it is the declaration of one person's initiative, not the result of an agreement between two parties, like a compact or a contract. This is beyond doubt one of the main reasons why the LXX rendered \he{b*:riyt} by \gr{διαθήκη} in the `covenants' of God; it was God alone who set the conditions; hence \textit{covenant} (see \hyperref[COVENANT-OED]{OED definition \#7}) can be used to translate \gr{διαθήκη} only when that is kept in mind. So \gr{διαθήκη} acquires a meaning in LXX which cannot be paralleled with certainty in extra-Biblical sources, namely `decree', `declaration of purpose', `set of regulations', etc. Our literature, which is very strongly influenced by LXX in this area, seems as a rule to have understood the word in these senses...
							\item[3] The meaning \textit{compact, contract} seems firmly established from Greco-Roman times...it remains doubtful whether this meaning has influenced our literature here and there...
						\end{compactitem}
					\item[CGL 345b, PDF 371] \textbf{}
						\begin{compactitem}
							\item[1] \textbf{state, condition}
							\item[2] document disposing of a deceased person's estate, \textbf{will}
							\item[3] deposits (apparently referring to sacred oracles, entrusted to the council of the Areopagus)
							\item[4] \textbf{agreement, covenant} (between individuals)
							\item[5] \textbf{decree} or \textbf{covenant} (of God) (New Testament)
						\end{compactitem}
					\item[LSJ 394b, PDF 465] \textbf{}
						\begin{compactitem}
							\item[I] \textit{disposition} of property by will, \textit{testament}
							\item[II1] mystic \textit{deposits} on which the common weal depended
							\item[II2] name of an eyesalve, because the recipe was \textit{deposited} in a temple
							\item[III] \textit{compact, covenant} \textit{(This entry includes references to LXX and the NT.)}
							\item[IV] = \gr{διάθεσις}
						\end{compactitem}
				\end{compactdesc}
				
				\begin{compactdesc}
					\item[Some NT uses] \textbf{}
						\begin{compactdesc}
							\item[Matthew 26:28] \gr{τοῦτο γάρ ἐστιν τὸ αἷμά μου τῆς διαθήκης τὸ περὶ πολλῶν ἐκχυννόμενον εἰς ἄφεσιν ἁμαρτιῶν·}
							\item[Mark 14:24] \gr{καὶ εἶπεν αὐτοῖς· Τοῦτό ἐστιν τὸ αἷμά μου τῆς διαθήκης τὸ ἐκχυννόμενον ὑπὲρ πολλῶν.}
							\item[Luke 22:20] \gr{καὶ τὸ ποτήριον ὡσαύτως μετὰ τὸ δειπνῆσαι, λέγων· Τοῦτο τὸ ποτήριον ἡ καινὴ διαθήκη ἐν τῷ αἵματί μου, τὸ ὑπὲρ ὑμῶν ἐκχυννόμενον.}
							\item[1 Corinthians 11:25] \gr{ὡσαύτως καὶ τὸ ποτήριον μετὰ τὸ δειπνῆσαι, λέγων· Τοῦτο τὸ ποτήριον ἡ καινὴ διαθήκη ἐστὶν ἐν τῷ ἐμῷ αἵματι· τοῦτο ποιεῖτε, ὁσάκις ἐὰν πίνητε, εἰς τὴν ἐμὴν ἀνάμνησιν.}
						\end{compactdesc}
				\end{compactdesc}
		
	% -----
	\subsubsection{Book of Mormon} % *BOM
	% -----
		\label{COVENANT-BOM}
	
		% \begin{tabular}{ll}
			% Total occurrences in the BOM & 
		% \end{tabular}
		
		\begin{compactdesc}
			\item[{\hyperref[2NEPHI-30-2]{2 Nephi 30:2}}]
		\end{compactdesc}
		
	% -----
	\subsubsection{Doctrine and Covenants} % *DC
	% -----
		\label{COVENANT-DC}
	
		% \begin{tabular}{ll}
			% Total occurrences in the DC & 
		% \end{tabular}
		
		\begin{compactdesc}
			\item[{\hyperref[DC22-1]{DC 22:1}}] BEHOLD, I say unto you that all old covenants have I caused to be done away in this thing; and this is a new and an everlasting covenant, even that which was from the beginning.
			\item[{\hyperref[DC40-1]{DC 40:1}}]
			\item[{\hyperref[DC54-4]{DC 54:4--6}}] And as the covenant which they made unto me has been broken, even so it has become void and of none effect. 5 And wo to him by whom this offense cometh, for it had been better for him that he had been drowned in the depth of the sea. 6 But blessed are they who have kept the covenant and observed the commandment, for they shall obtain mercy.
			\item[{\hyperref[DC66-2]{DC 66:2}}]
		\end{compactdesc}
		
	% % -----
	% \subsubsection{Pearl of Great Price} % *PGP
	% % -----
		% \label{COVENANT-PGP}
	
		% \begin{tabular}{ll}
			% Total occurrences in the PGP & 
		% \end{tabular}
		
		% \begin{compactdesc}
			% \item[]
		% \end{compactdesc}
		
% ----------------------
\subsection{Key Phrases} % *KEYPHRASES *PHRASES
% ----------------------
	\label{COVENANT-PHRASES}

	\begin{compactdesc}
		\item[new covenant] \textbf{16} | Jeremiah 31:31; Hebrews 8:8, 13; 12:24; DC 22:1; DC 76:69; DC 84:57; DC 107:19; DC 131:2; DC 132:4, 6, 19, 26, 27, 41, 42
			% \begin{compactdesc}
				% \item[Bible anchor verses] \phantom{.}
					% \begin{compactdesc}
						% \item[Romans 5:5] \gr{}
					% \end{compactdesc}
				% \item[Commentary] ---
			% \end{compactdesc}
		\item[everlasting covenant] \textbf{35} | 
			\begin{compactdesc}
				% \item[Bible anchor verses] \phantom{.}
					% \begin{compactdesc}
						% \item[Romans 5:5] \gr{}
					% \end{compactdesc}
				\item[Commentary] The ``everlasting covenant'' and the ``fulness of the gospel'' are equated in \hyperref[DC66-2]{DC 66:2}.
			\end{compactdesc}
	\end{compactdesc}
	
% % ----------------------
% \subsection{Bible Dictionaries} % *DICTIONARIES
% % ----------------------
	% \label{COVENANT-BD}

% % ----------------------
% \subsection{Restoration Usage} % *RESTORATION
% % ----------------------
	% \label{COVENANT-RESTORATION}

	% % -----
	% \subsubsection{Joseph Smith} % *JS
	% % -----
		% \label{COVENANT-JS}
	
		% \begin{compactdesc}
			% \item[{\hyperref[KEY]{Date/Source}}]
		% \end{compactdesc}
	
	% % -----
	% \subsubsection{Brigham Young} % *BY
	% % -----
		% \label{COVENANT-BY}
	
		% \begin{compactdesc}
			% \item[{\hyperref[KEY]{Date/Source}}]
		% \end{compactdesc}
	
% ----------------------
\subsection{Commentary and Studies} % *COMMENTARY *STUDIES
% ----------------------
	\label{COVENANT-COMMENTARY}

	% -----
	% \subsubsection{Key Points}
	% -----
		% \label{COVENANT-KEYS}
	
		% \begin{compactitem}
			% \item 
		% \end{compactitem}
		
	% -----
	\subsubsection{Why in general are covenants necessary?}
	% -----
		\label{COVENANT-necessity}
		
		For the view I most prefer on this issue, or for the beginnings of it at least, see \href{https://climbingtherainbows.substack.com/p/naughty-nellys-home-4-teh-holidayz-tour-2024-nels-nelson-on-covenants}{this essay}, which draws from the work of Nels Nelson.
		
		Some applicable scriptures:
		
		\begin{compactitem}
			\item \hyperref[DC18-25]{DC 18:25}, you have to ``know the name,'' or be consciously aware of having chosen to follow Christ
		\end{compactitem}
		
	% -----
	\subsubsection{The ``new and everlasting'' covenant}
	% -----
		\label{COVENANT-everlasting}
		
		See JSP-J1:72--73 (18 October 1835) for a journal entry talking about blessing children with the blessings of the ``new and everlasting covenant.'' See also 1 November 1835; \hyperref[DC22-1]{DC 22:1} (this verse is important because it makes it clear that the ``new'' covenant is really the \textit{original} covenant, from the very beginning); 
		
	% -----
	\subsubsection{More reference information about Hebrew and Greek words}
	% -----
		\label{COVENANT-reference}
		
		% --
		\paragraph{NIDNTT I:365--372}
		% --
			\label{COVENANT-NIDNTT}
		
			\textit{(Note that the editor of NIDNTT makes several clarifications and even challenges the author's interpretation at some points, so this entry may be a bit idiosyncratic.)}
		
			%
			\subparagraph{Classical sources}
			%
			
				The term diatheke occurs from Democritus and Aristoph. onwards in the sense of a will or testament. It is not thought to be derived from the act. diatithemi, distribute, allocate, regulate, but only from the mid. diatithemai, control persons and things (Xen.), and especially dispose of by will (so private legal documents among papyri). It denotes, therefore, an irrevocable decision, which cannot be cancelled by anyone. A prerequisite of its effectiveness before the law is the death of the disposer. Rence diatheke must be clearly distinguished from syntheke, an agreement. In the latter two partners engaged in common activity accept reciprocal obligations. diatheke is found only once with this meaning (Aristoph. Birds, 1,440). Elsewhere it always means a one-sided action. 

			%
			\subparagraph{Old Testament}
			%
			
				1. In the LXX diatheke is the commonest rendering (270 times) for Heb. \he{b*:riyt}, covenant. This is the common OT word for a wide variety of agreements. (a) They include a covenant between two friends (1 Sam. 18: 3), which however was regarded as having legal force (1 Sam. 20: 8); a covenant between two rulers fixing their spheres of interest (Gen. 21: 22 ff.; 26: 26 ff.; 1 Ki. 5 :12, MT 5: 26), or terms of peace (1 Ki. 20: 34); the covenant between two kings which, of course, included their peoples. Two tribes could also make a covenant (Jos. 15: 9). It is even used for a covenant between Israel and its slaves (Jer. 34: 8 ff.). ([Tr.] The I natural meaning of the Heb. is that it was a covenant between the king and the slave owners.) 

				(b) The situation was somewhat different when the king made a covenant with his subjects (2 Ki. 11 :4; [Tr.] this was probably a private agreement between Jehoiada and the royal guard); or when by the covenant a man became the subject of the other (1 Sam. 11 :1). The covenant made by Abner with David (2 Sam. 3 :12 f.) was intended to make all Israel (v. 21) subject to David, not merely Abner. 

				(c) Yahweh's covenants with Noah (Gen. 6:18), Abraham (2 Ki. 13:23), or David (Jer. 33: 21) are similar. Here the covenant extends explicitly to their descendants (Gen. 9:8 ff.; 15:18; cf. 2 Sam. 7:12-16), and becomes a covenant with Israel (Exod. 6: 4 f.). But in Jer. 50: 5 the covenant can also be interpreted as a covenant of Israel with its God. Ezek. 16: 8 speaks of a covenant with Jerusalem. Linguistically the LXX might have used syntheke in the examples under (a) (cf. CL above). But in the latter cases berit is used to denote the one's disposition of the many (Gen. 9 :12; 17: 7 f.; Jer. 31: 31-34, etc.). Hence the LXX translation diatheke has its full force. On the other hand, there are exceptional usages like Yahweh's covenant with day and night (Jer. 33: 25), and Jerusalem's with death (Isa. 28: 1418). 

				(d) Earlier scholars used to distinguish between secular and religious covenants but the variety of covenant conceptions, all using the same terminology, makes this impossible. Moreover, the most secular of covenants presupposes that it has been witnessed by God (Gen. 31 :44-50; 1 Sam. 23:18), while in the most religious covenants Yahweh remains comparable with a monarch making a covenant. For the aT a separation between religious and secular is unthinkable. 

				2. This situation has made many ask whether there was not some original and unambiguous idea of the covenant, from which the others developed. Modern research has produced four answers. (i) J. Pedersen considered that originally the two sides in the covenant stood on an equality. (ii) J. Begrich, on the other hand, maintained that the covenant started as something the stronger granted to the weaker. (iii) G. E. Mendenhall found the religio-historical background in the Hittite suzerainty treaties, and deduced that the treaty element was the most important in the covenant. (iv) M. Noth has drawn attention to the most recent textual finds, which show that the covenant was mediated by a third party between the two sides. It is impossible to prove that anyone of these elements was original, or that any came later. God's covenant and the people's covenant are mentioned in Exod. 34; the individual and the royal covenant in 2 Sam. 3. We cannot at present decide on any original form for the covenant. 

				3. Since all these concepts were expressed by the one term berit, the Jews must have felt some unity behind them. Though neither the relationship of the covenant partners nor the contents of the covenant agreement provided this, it has been shown that there was a common pattern in the making of a covenant (Baltzer, see bibliography). There were 6 vital elements in such a ceremony: (i) "the preamble mentioning the names of the partners; (ii) a preliminary history of the relationship of those entering the covenant; (iii) a basic declaration about the future relationship of the partners; (iv) details of the new relationship; (v) an invocation of the respective gods worshipped by both sides to act as witnesses; (vi) a pronouncement of curse and blessing. Baltzer traced this outline right down to the Qumran documents, and even into and beyond the NT. This shows that at least outwardly the covenant conception is as fundamental an element in biblical thought as it is in the national existence of Israel. 

				4. When we approach them from the standpoint of the covenant, we see how very diverse elements in OT tradition belong together. The preamble is suggestive of the way in which Yahweh introduces himself (cf. Exod. 20:2, which, as in the covenant formula, is followed by a survey of past events). This seems to be the source of Israel's unique interest in its history, which is seen to be Yahweh's history with it, and hence covenant history and so salvation history. The third section reflects the thought which found its clearest expression in Isa. as proclamation of salvation, the promise of diville faithfulness and peace (Isa. 54: 10; 55: 3). Since this promise cannot be revoked because of the disobedience of the people, God has proclaimed himself bound and has renounced all arbitrariness. "The basis is thus laid for the Gospel" (G. Quell, TDNT II 123). When we compare these covenant promises and those in Jer. 31 with (e.g.) the formula of the Sinai covenant we realize that the details mentioned in 3 (iv) above are not to be understood as a condition to be first fulfilled which would be followed by a reward for achievement. This impression might be given by passages like Lev. 26:15 f. and Deut. 31 :20. They are rather the regulations for the new life that is made possible by God's covenant (cf. Gen. 9:9; 15:13-16; Jer. 31 :31). As in Deut. 26, Israel's reflection on its past history and the law belong together as instruction for its new life. 

				5. The relationship of the partners in the covenant is expressed by /:zeseg, God's covenant loyalty (RSV steadfast love; NEB love, loyalty, constancy). 1 Sam. 20: 8 understands it as protective action. Man's remembrance of the covenant expresses itself in action (Ps. 103: 17 f.). Both partners Yahweh and the covenant people, represented only by an individual face one another in the beril. They are thus in an active and real partnership (Pedersen), and so they both share in the covenant meal (Gen. 31: 54; Exod. 24:9 ff.). It goes without saying that this strengthens the -+ fellowship of those involved. This does not mean to say that the covenant and its renewal were purely cultic acts. It can hardly be denied that it had a place in the cultus (Ps. 50). This seems to have been especially the case in Shechem (Jos. 8: 3035; 23: 24). The relationship to Baal-berith or El-berith, formerly worshipped there, has never been satisfactorily explained (Jdg. 8: 33; 9: 4, 46). Nevertheless, the covenant belongs in the first place to the daily life of the tribes of the covenant people, which it seems to have steadily separated from the cultus. This may be seen in Deut. with its covenant theology and secularization of life. The cultus could not guarantee the continuance of an intact, everlasting covenant. Yahweh alone, as founder of this covenant, could guarantee its continuance and with it the cultus in its true meaning. Only he could renew the covenant broken by human disobedience. 

				This explains why in many prophets the thought of the covenant virtually vanishes in favour of the call to obedience. The early parts of Isa. avoided obviously deliberately the concept of covenant almost entirely, the only exceptions being 24:5; 33:8.and the covenant with death mentioned earlier (28:15, 18). Hosea probably felt the danger of a one-sided appeal to the legal character of the covenant, and so he preferred to describe God's relationship to his people by the picture of marriage. The berit is mentioned only as broken in 6: 7; 8: 1; in 10: 4; 12:1(2) it refers to human agreements. Jeremiah stressed his expectation of a completely new covenant in the place of the one that had been broken (31: 31-34). Whatever continuity there may be is created solely by Yahweh. It is completely other in that the direct working of God on human hearts through his Spirit takes the place of the demand for obedience. This shows how deep the break is, and prepares the way for the NT. 

				([Ed.] There is, however, an essential continuity, for Jeremiah does not speak of the abolition of the law. Rather, he declares: "But this is the covenant which I will make with the house of Israel after those, says the Lord: I will put my law within them, and I will write it upon their hearts; and I will be their God, and they shall be my people" (Jer. 31: 33). There is here the same divine upholding of the covenant that is found elsewhere in the OT (see 4, above). The covenant relationship is summed up in the promise, "I will be your God, and you shall be my people" (cf. Exod. 6:7; Lev. 26:12; Jer. 7:23; 11 :4; 30:22; 32:38; 2 Cor. 6:16; Heb. 8:10; Rev. 21: 3). For Jeremiah, this relationship is to be realized not by setting aside the law but by a more personal application of it.) 

				6. The standard expression in the OT, kiiral beril, to cut, i.e. make, a covenant, is probably to be explained by Gen. 15:9 f., 17 and Jer. 34:18. ([Tr.] It is not the only expression used, and it does not seem to have been essential for the making of a covenant.) It is expressed in the LXX by diatheken diatithemai, to establish a diatheke. The Gk. understanding of the term diatheke suggests that the covenant was no longer an agreement between two parties with equal rights. It came about as an exclusively divine action, which men can only accept in the form in which it is given to them. However, the LXX is supported by the Heb. expression heqim berit, which occurs frequently after the exile, and means "to establish a covenant". The apocryphal books carry this further. Philo worked out an allegorical exegesis of the thought of a will. The rabbis confined the covenant to the concept of an absolute law. 

				The Qumran community attributed great importance to the covenant. They considered that the promise of the new covenant (Jer. 31: 31-34) had been fulfilled in their midst and called themselves "the new covenant in the lane! of Damascus" (CD 6:19). While the Sinaitic covenant was for the whole people, the Qumran covenanters considered themselves to be its holy remnant, the pure, eschatological community of the age of salvation. This explains their strict rules for receiving new members and their commitment to the "laws of the covenant (i.e. the Law as it was expounded by the community) which characterized the community. The sect's eschatological thinki.ng on the covenant distinguished them especially from the rabbis, who equated the covenant with circumcision and saw its preservation guaranteed by a rigorous keeping of the Mosaic Law, and who expected the fulfilment of Jer. 31: 31-34 only in the future. ([Tr.] In fact, probably owing to polemics against Qumran and Jewish Christianity, the new covenant plays virtually no role in Talmudic Jud.) 

			%
			\subparagraph{New Testament}
			%
			
				It is noteworthy that while covenant is found almost 300 times in the OT, it occurs only 33 times in the NT. Almost half of these instances come in quotations from the OT, and another 5 clearly look back to OT statements. The few independent cases are almost exclusively in Heb., quite rarely in Paul (never in the Pastorals) and Acts, and never in the Johannine writings. 

				In accord with this, diatithemai (only mid.), order, enact, dispose, is found only twice in the Lukan writings and four times in Heb., each time clearly looking back to the ~T. In the 2 cases where the standard Gk. meaning of "will" plays a part (Heb. 9: 16; Gal. 3: 15) it is easy to recognize the intention of helping the Greeks to understand the OT concept. Paul makes this explicit. 

				These statistics are, however, misleading. The covenant question in the NT cannot be answered solely from the passages where the word is used. It involves a whole complex of theological ideas including covenant terminology. There are 3 groups of problems: (1) the question ofthe Lord's Supper; (2) the Pauline question about the relationship of the Christian church to Israel as the people of God; (3) the covenant theology in Heb. 

				1. In all 4 NT passages dealing with the Lord's Supper diatheke plays an important part: 1 Cor. 11: 25; Mk. 14: 24; Matt. 26: 28; Lk. 22: 20 (omitted in some MSS and RSV and NEB). Each time it is in connection with the cup, and only there. Paul and Lk. add the adj. "new". This gives a new stress, which hides to some extent the link with the blood, in contrast with the original formula preserved in Matt., Mk., "the blood of the covenant". The use of this formula together with "shed" or "poured out" makes clear beyond question the reference to the covenant blood of the OT (cf. Exod. 24: 5-8) and with it the covenant which Yahweh made with Israel. This means that, in the Christian kerygma and witness, the work of Jesus was, according to his own word, a taking up and fulfilling of the covenant statements of the ~T. Schniewind is, therefore, fully justified in pointing out that "Jer. 31: 31-34 can be heard in every variation of the words over the cup" and that the prophetic promise finds fulfilment here. 

				If the term covenant (diatheke) does not appear as often as one might expect, the reason is that the underlying thought has been taken over in the sayings about the ~ kingdom of God. Linguistically we can see this perhaps most clearly in Lk. 22: 29 in the phrase diatithemai ... basileian, appoint a kingdom, which exactly expresses the formula diatithemai diatheken (see above OT 6). The new covenant and the kingdom of God are correlated concepts. 

				In the OT the common meal followed the completion of the covenant by the sprinkling of the people with blood after its pledge (Exod. 24:8 if.), because the covenant had been made with a people that was already united. In the NT such a single, united people cannot be presupposed. The new people of God is derived rather from that which until then was no people (Rom. 9: 25; 1 Pet. 2: 10, quoting Hos. 2: 23 (25». It is composed of the "many" who have come to share the atoning power of Jesus' blood through his word and spirit (Matt. 20: 28; Mk. 10: 45). The fellowship of the early Christians based on the word found expression in the common meal (1 Cor. IO:I6b, 17; Acts 2:46), which has its counterpart in the first part of the Lord's Supper. 

				As in the old covenant curse and blessing played a part, so there is also in the new covenant an indication of the curse (1 Cor. 11 :27-32; cf. 12: 3). Just as one could keep the old covenant only if one knew its ordinances, so in the Lord's Supper repetition is expressly intended to produce remembrance (anamnesis, I Cor. 11: 25). Not long after this entire Christian writings like the Didache were moulded by the theology of the Lord's Supper and show the structure of the covenant ritual. It is only a development of a connection that is already obvious here. 

				2. It is not easy to formulate Paul's attitude to the old people of God. 

				(a) Paul himself can use only the term "mystery" (---+ mysterion, secret) for the Jewish rejection of Jesus and their enmity to the gospel and the church (Rom. 11 : 25). He does so in "great sorrow and unceasing anguish" (Rom. 9: 2). He was so deeply moved by this subject, that his delineation of the righteousness of God given through Christ (Rom. 1-8) of necessity passed over to a presentation of the inexplicable hardening of Israel, which nevertheless is to lead to salvation in God's will (Rom. 9-11, especially 11 : 11-16). This account reaches its climax in praise of the unfathomable depths of the divine action (Rom. 11: 33-36). In it Paul appeals to the fact that Israel's rejection and unfaithfulness cannot cancel God's covenant. He also knows that there is a fellowship between Israel and the church which cannot be removed, because it is based on the unity of the One who calls. In other words, God's people from now on includes the church and Israel. 

				(b) In the light of this, Paul considers the fact that a part of Israel has been hardened (Rom. 11: 25). The expression "a hardening has come upon Israel" (to Israel gegonen) suggests that it has happened through God's will. God promised Israel in his covenant that sin would be completely blotted out and that all obstacles between God and man would be removed without trace. This is the force of the quotation from Isa. 59: 20 f. in Rom. 11: 26 f. It is, therefore, obvious that the fulfilment of this covenant promise is still to be awaited, so long as the obstacle of the hardening still exists. Since God's covenant has not reached complete fulfilment or has been cancelled or repented of by him (Rom. 11: 29), the promise will be fulfilled with regard to the remnant of Israel. 

				(c) To make this clearer still, in Gal. 3: 15 Paul examined the Gk. concept of a will that could not be annulled. If a human will cannot be annulled, how much less can God's covenant with Israel (3: 17), through which it has such incomparable privileges (Rom. 9: 4 f.), of which Paul the Christian can boast again and again (Rom. 11:1; 2 Cor. 11: 22; Phil. 3: 5)? 

				If God's covenant with Abraham, which consists of the promise, has come into force (Gal. 3:16), something which came in later as an addition to it (Rom. 5:20), the law (Gal. 3 :17), cannot narrow this promise or annul it. In the opinion of the present writer, Paul could hardly mean the law which according to OT conceptions belonged to the covenant by its very nature, and thus must be regarded as an ordinnance (see above OT 4). Hence it is no more in contradiction with the covenant of grace as the OT is with the NT. Paul is thinking rather of the law as seen by late Jud., and the rabbis and understood as something in its own right and absolute (see above OT 6). This mass of inexplicable regulations, which simply have to be fulfilled, has been annulled by Christ (Rom. 10: 4), because it was only a paidagogos, custodian (RSV), tutor (NEB), "until Christ came" (Gal. 3: 24). The law of God itself is not evil, but "holy and just and good" (Rom. 7: 12). It belongs to the Sinaitic covenant (Gal. 4:24), not to the Abrahamic covenant (Gal. 4:23, the picture of the two women representing two covenants). Hence it cannot bring salvation. The one who clings to it remains a slave child (Gal. 4: 25). Christians, on the other hand, are children of the promise (Gal. 4:26 ff.) and hence free, and belong to Christ without any distinctions (Gal. 3: 26 ff.). 

				([Ed.] The identification of the law here with the elaborations of late Judaism is not the most obvious interpretation of the references cited. The argument about circumcision and keeping the law in Rom. and Gal. does not suggest that Paul is arguing against a particular rabbinic interpretation. Rather, he is arguing that man cannot use the ordinances of God to establish his own righteousness. Man is saved by the grace of God through Christ's atoning death. Paul does not attack the ordinances of the law as such, but the assumption that man becomes righteous through conformity with certain tenets. Such an attitude of mind implies that it is in man's power to keep the entire law. It betrays a profound lack of awareness of the deep-rooted character of sin and its seriousness in the eyes of God. The Christian, Paul argues, must not relapse into a legalistic frame of mind, for righteousness is the free gift of God in Christ.) 

				3. We have been seeing how the covenant concept stands in its relationship to the -+ Law, the -+ Promise, and the people of God, but we have to turn to Heb. for a developed covenant theology (diatheke is mentioned 17 times). Since the highpriestly office of Christ is "the prevailing conception of Hebrews' Christology", it expresses the new covenant in a cultic setting, which reflects the author's intense concern for "purification, sanctification, perfection" and also sacrifice, atonement and blood (W. G. Kiimmel, Introduction to the New Testament, 1966, 277). He clearly contrasts the old and the new covenants. The new is the better covenant (cf. 7:22 kreitt6n, superior). Its guarantor and mediator is Christ (8:6; 12:24) through his death for the redemption from the sins of the old covenant (9: 5). It is against this background of blemish that the prophetic promise of Jer. 31: 31-34, twice quoted (8: 8-12; 10: 16 f.), becomes important and is developed. By the very promise of the new covenant the old has been declared obsolete by God himself, and so it is ready to vanish away (8: 13). This does not mean that it has been completely ruled out, but that it has been overtaken and fulfilled by the new. In Christ the pattern (typos, 8: 5) and "the true form" (eik6n, 10: 1) of the one covenant concept has become a reality. The new covenant is founded on better promises (8:6; 9:15), but like the old covenant it is a covenant in blood (10:29; 12:24; 13 : 20). It is not, however, the blood of a sacrificial animal, but the blood of the sacrificer himself, the high priest (9:13 ff.; 10:12 ff.). Hence the death of Christ was essential, for only through the death of the testator does his will become operative in law (9:16 f.; cf. above CL). 

				Seeing that the forgiveness of sins and renewal of heart promised in Jer. 31: 3134 has become a reality in Christ (Heb. 10:16 ff.), the old covenant is explained as a shadow of that to come. In the new it has been annulled in the sense that it has become "obsolete" and "superfluous". On the other hand, the Pauline tension over the problem of Israel's hardening has disappeared here. That which Paul could only proclaim in his praise of God's mystery (Rom. 11: 33-36) is presented in Heb. in an almost rationalistic and comprehensive manner. This can be explained probably only because the author of Heb. faces the old covenant more theoretically, while Paul, in spite of all, knows and confesses that he belongs to it. 

				([Ed.] Many scholars would dissent from the writer's conclusions here. It is not that the author of Heb. presents everything more rationalistically than Paul. Indeed, he refrains from speculating about the future of Israel and does not discuss the problems that Paul wrestled with. Instead, his warnings against the danger of rejecting Christ come like a refrain throughout his work (3:7-19; 4:1-13; 6:4-8; 10:26 ff., 38 f.; 12:1-29). Conscious of the judgments which befell the rebellious and hardened in Israel, he warns his readers against deliberate, conscious sin. For him there is a point beyond which a man cannot return. That point cannot be defined in advance. The safeguard against falling into condemnation is summed up in the quotation from Ps. 95:7 f., "Today, when you hear his voice, do not harden your hearts" (Heb. 3:7, 15; 4:7). 
				
		% --
		\paragraph{TDOT II:253--279}
		% --
			\label{COVENANT-TDOT}
			
			The original meaning of the Heb. \he{b*:riyt} (as well as of Akkadian \textit{riksu} and Hittite \textit{i\u{s}\b{h}iul}) is not ``agreement or settlement between two parties,'' as is commonly argued. \he{b*:riyt} implies first and foremost the notion of ``imposition,'' ``liability,'' or ``obligation,'' as might be learned from the ``bond'' etymology discussed above. Thus we find that the \he{b*:riyt} is commanded (\textit{tsivvah beritho}, ``he has commanded his covenant,'' \hyperref[PSALMS-111-9]{Psalms 111:9}, \hyperref[JUDGES-2-20]{Judges 2:20}), which certainly cannot be said about a mutual agreement. As will be shown below, \he{b*:riyt} is synonymous with law and commandment (cf., e.g., Deuteronomy 4:13; 33:9; Isaiah 24:5; Psalms 50:16; 103:18), and the covenant at Sinai in \hyperref[EXODUS-24]{Exodus 24} is in its essence an imposition of laws and obligations upon the people (\hyperref[EXODUS-24-3]{verses 3--8}).
			
			